% =====================================================================
%  Introduzione  (di Michela Mollia) — scansione p-01 .. p-04
%  Testo trascritto integralmente dalle immagini (de-ombreggiate).
% =====================================================================
\chapter*{Introduzione}
\addcontentsline{toc}{chapter}{Introduzione}
\markboth{Introduzione}{Introduzione}

Un sistema musicale è un'organizzazione di relazioni (intervalli) tra altezze
e inevitabilmente tali relazioni vengono rappresentate da numeri. Il calcolo
degli intervalli non è certo materia nuova, ma nemmeno estranea ad un certo
tipo di pensiero compositivo contemporaneo, come vedremo più avanti.

In effetti esso affonda le sue radici nell'antichità più remota. Pitagora da
Samo nel IV secolo a.C. ne aveva stabilito alcune delle regole fondamentali,
ma certamente civiltà musicali diverse, come ad esempio quella cinese, ne
conoscevano i principi: una corda viene suddivisa in parti via via sempre più
piccole, prima la metà, poi un terzo, un quarto\dots\ la scoperta di altezze
diverse le cui relazioni con la nota fondamentale sono le stesse con cui gli
armonici si mettono in rapporto l'uno con l'altro rispetto al fondamentale: ma
con una piccola grande differenza, inversamente proporzionali. Questa legge
semplice non è altro che il principio degli armonici riscontrato e formulato
solamente nel XVII secolo. È comunque sorprendente notarne l'antica intuizione.

L'estrema semplicità di questo generatore, che ancora non sa di musica, di
fantasia, di amore o di guerra, è però perfetta nella sua regolare disposizione
di leggi. Vale la pena ricordare come tale equilibrio venne poi ricercato nel
movimento degli astri, nelle forme geometriche più semplici, nel modo come la
natura ordina e costituisce i suoi organismi, dai più elementari ai più
complessi. A questo punto sarebbe fin troppo facile cantare non originali lodi
sulla meravigliosa unitarietà dell'universo. La nostra dimensione ha perduto,
purtroppo o no, questo legame, la consapevolezza di far parte di un unico tutto.
Abbiamo altre cose per le mani.

I nostri strumenti musicali sono molto più complicati, ma se togliamo loro le
accurate e sofisticate sovrastrutture di legno o metallo, troveremo la familiare
corda vibrante o l'inafferabilità della colonna d'aria. Questi oggetti sono per
lo più legati indissolubilmente a leggi e modelli fisico-musicali che hanno
provocato e permesso sia la loro realizzazione sia opere che di questi strumenti
facessero uso.

Quando anni fa il compositore cominciò a modellarsi la materia sonora
nell'ambito delle esperienze della musica concreta, era in un certo senso
ancora indifferente alla sua quantizzazione interna, lavorando più su aggregati
timbrici, più o meno evocativi, che su grandezze sonore su cui fosse oggetto una
rete di relazioni, se pure arbitraria, che costituisse un sistema. La dimensione
sonoro-compositiva muta, direi radicalmente, nel momento in cui si lavora con
generatori elettronici analogici e ancora di più, recentemente, con elaboratori
digitali che, tramite adeguati sistemi di conversione numerico/analogica, sono
in grado di sintetizzare qualsiasi forma d'onda.

Se a questo aggiungiamo la possibilità di creare di volta in volta dei sistemi
di intervalli che soddisfino le necessità espressive e compositive del momento
(aprendo magari la strada ad una grammatica degli intervalli) si può intravedere
un universo sonoro molto più sfaccettato e versatile di quanto il temperamento
equabile non abbia permesso.

In effetti, se a questo punto facciamo un'analisi musicologica dei sistemi di
intonazione che si sono imposti nella pratica musicale occidentale, vediamo che
soprattutto tre hanno avuto una rilevanza storica. Fino al XVI secolo ci si
avvale del sistema pitagorico, che considera solamente i numeri 1, 2, 3 e 4
(fondamentale, ottava, quinta e doppia ottava). Questa serie di numeri naturali
viene chiamata appunto \emph{Tetractys} pitagorica. Essa costituiva il nucleo
dal quale si faceva risalire la generazione e il susseguente ordinamento
armonico del mondo. I pitagorici concepivano il mondo come musicalmente
costituito e fissavano questo concetto con il termine \enquote{cosmo}.

Con il formarsi dell'armonia moderna questo sistema si rivelò essere di
difficile impiego, in quanto non facilitava un'organizzazione basata su accordi
(la consonanza perfetta non poteva essere di soli due suoni). Al contrario,
favoriva un contrappunto di linee indipendenti. Ed ecco perché la polifonia
primitiva, come pure quella del Medio Evo, è una polifonia di contrappunto,
guidata dal procedere melodico delle parti e non dai rapporti contenuti negli
accordi.

Altra caratteristica del sistema pitagorico è quella di avere il semitono
considerevolmente più piccolo di quello temperato, mentre tutti gli altri
intervalli risultano allargati. Questo provoca una tendenza molto forte di
alcuni di questi a risolvere su altri più semplici. Il sistema pitagorico è
quindi essenzialmente dinamico, accentua le differenze tra gli intervalli e
aumenta le dissonanze, che devono perentoriamente risolvere.

Il sistema cosiddetto zarliniano opera un allargamento armonico da 2, 3 suoni,
giungendo così al fattore 5 (terza maggiore). Questo intervallo rende l'accordo
più stabile, passando dal rapporto $81/64$ a quello più consonante di $5/4$.
L'armonia del XVI secolo liquida progressivamente il contrappunto e, smussando
le attrazioni intervallari pitagoriche, forma un tappeto di accordi consonanti
che fanno risaltare la melodia. Possiamo perciò dire che il sistema zarliniano
è, in tal senso, un sistema statico: la gerarchia dei cicli delle quinte finisce
per scomparire, la struttura cessa di essere do-fa-sol-do per diventare il
do-mi-sol-do.

Differentemente dai precedenti, il temperamento equabile non è né statico né
dinamico. Esso opera dei continui compromessi tra i due sistemi: un sistema
neutro nel quale, operando mentalmente continue correzioni degli intervalli, ci
muoviamo melodicamente verso il pitagorismo mentre per l'armonia ci rivolgiamo
al sistema zarliniano. Esso non impone nulla e permette tutto. D'altra parte
questa ambiguità ha permesso di suonare in tutte le tonalità, di modulare
rapidamente e fare uso dell'enarmonia in senso moderno. A quanto detto, se pur
brevemente, si può aggiungere di conseguenza come il passaggio nella storia
della musica da un sistema di intonazione ad un altro non sia un fatto
accidentale, ma uno degli elementi essenziali, la matrice determinante
dell'evoluzione del linguaggio musicale.

D'altra parte un compositore attento, come era senza dubbio Arnold Schoenberg,
nel suo saggio \emph{Problems of Harmony} del 1934 osservava come noi stiamo
provando ad esprimere armonie del 7, 9 e 11 in un sistema e con strumenti di
notazione formati soltanto per quelle del 3 e del 5, sottolineando sia
l'incapacità del temperamento equabile di reggere la spinta evolutiva di un
linguaggio musicale che stava cambiando, sia la contraddizione tra un operare,
diciamo così, avanzato, che però si serviva di un sistema nato per tutt'altre
esigenze.

A questo punto il saper trattare ed esporre, essenzialmente e non con la pretesa
di teorici, lo studio e la pratica riguardante il calcolo degli intervalli ci
sembra molto proficuo e in definitiva affascinante. Perlomeno per quella libertà
di definizione che oggi ci è permessa da mezzi come gli elaboratori, che
garantiscono la precisa realizzazione di qualsiasi grandezza voluta. Conoscendo
le caratteristiche dei vari sistemi di intonazione sceglieremo quelli, o parte
di quelli, più adatti alle nostre esigenze compositive.

Inoltre, dato che nel campo della musica elettronica non esistono suoni
preformati, è ancora più evidente come chi opera in questo ambito debba, per
forza di cose, predisporsi tutti i livelli della composizione. E dato che ai
calcolatori si parla con i numeri (almeno nella maggior parte dei casi) si può
ritrovare come estensione quella antica concezione di numero e suono in cui uno
è espressione dell'altro e dove qualità e quantità si trovano per l'unica volta
unite. In altre parole, quando parliamo di \enquote{ottava} ci riferiamo
inequivocabilmente a quel certo intervallo, a quella precisa qualità
intervallare, ma, nello stesso tempo, a quella precisa quantità, misurazione se
vogliamo, di corda o colonna d'aria.

Chiediamo ad un musicista di misurare una corda e di indicare la sua esatta metà
senza l'uso di un sistema di misurazione: questo calcolo verrà effettuato nella
maniera più empirica ed efficace, andandosi a cercare cioè l'ottava. Si potrebbe
dire perciò con Curt Sachs: \enquote{Il suono non ha relazione diretta con altre
forme di percezione, avendo rispondenza solo con gli aspetti più astratti della
realtà: il numero e la misura}.

Il calcolo degli intervalli rappresenta in musica ciò che il problema della
quadratura del cerchio è in geometria: un calcolo che si avvicina sempre più al
segno di infinito. D'altra parte, se è vero che la serie dei numeri è
praticamente illimitata, possiamo pensare ad una illimitata serie di intervalli?
Si può dire che il nostro orecchio pone delle condizioni di discriminabilità: la
minima differenza di \cents{6} costituisce una soglia più o meno generalizzabile
per il riconoscimento degli intervalli. Certamente, comunque, con un computer
potremmo conoscere e cogliere quelle minime differenze, giungendo, perché no, ad
un tipo diverso di espressività, che sembrava estranea ad uno strumento da molti
considerato una macchina.

Scopo quindi di questa trattazione è quello di fornire gli elementi e principi
essenziali di calcolo degli intervalli. Questa conoscenza può essere di grande
importanza sia per un'analisi più appropriata dei sistemi di intonazione già
messi in pratica, che come punto di partenza per la creazione di nuovi.

Attraverso queste pagine è possibile seguire il cammino della musica così come
\enquote{suonava} tanto tempo fa, che è sicuramente molto diverso da quello di
oggi. Il concetto di intervallo, quindi, non è da leggersi come irrevocabilmente
legato alle \enquote{note}, bensì va inteso nel senso più allargato di
proporzione, distanza, emozione.

\firma{Michela Mollia}
